% lecture24.tex — Gauss-Bonnet Theorem

\section{Gauss--Bonnet Theorem}

Let $M$ be a compact hypersurface (i.e.\ $m$-dimensional submanifold) in $\R^{m+1}$.

\begin{definition}{Gauss Map}{gauss-map}
    The \emph{Gauss map}
    \[
        g \colon M \to S^m, \qquad x \mapsto \vec{n}(x),
    \]
    sends each point $x \in M$ to the outward-pointing unit normal vector $\vec{n}(x) \in S^m$.
\end{definition}

\begin{definition}{Volume Form}{volume-form-Gauss}
    The \emph{volume form} $\mathrm{vol}_M$ on $M$ is the $m$-form on $M$ which evaluates to $\frac{1}{m!}$ on each positively oriented orthonormal basis of $T_x M$. Define $\mathrm{vol}_{S^m}$ on $S^m$ in the same way.
\end{definition}

Since they are top-dimensional forms, at each point $x \in M$, $g^* \mathrm{vol}_{S^m}$ must be a scalar multiple of $\mathrm{vol}_M$.

\begin{definition}{Gaussian Curvature}{gaussian-curvature}
    The \emph{Gaussian curvature} of $M$ is the function $\kappa \colon M \to \R$ such that
    \[
        g^* \mathrm{vol}_{S^m} = \kappa\, \mathrm{vol}_M.
    \]
    It measures how much $g$ distorts the volume (and orientation).
\end{definition}

\begin{example}{Hypersurfaces in $\R^3$}{Gauss-curvature-R3}
    For a convex surface (e.g.\ a sphere or ellipsoid), nearby normals on $M$ map to nearby normals on $S^2$ in an orientation-preserving way, so $\kappa > 0$. For a saddle surface, the Gauss map reverses orientation locally, so $\kappa < 0$.
\end{example}

\subsection{The Gauss--Bonnet Theorem}

The Gauss--Bonnet theorem relates the Gaussian curvature (geometric, local) with the Euler characteristic (global, topological).

\begin{theorem}{Gauss--Bonnet}{gauss-bonnet}
    If $M$ is a compact, even-dimensional hypersurface in $\R^{m+1}$, then
    \[
        \int_M \kappa = \tfrac{1}{2}\, \gamma_m\, \chi(M),
    \]
    where
    \[
        \gamma_m := \int_{S^m} \mathrm{vol}_{S^m} = \frac{2 \cdot (2\pi)^{m/2}}{1 \cdot 3 \cdot 5 \cdots (m-1)}
    \]
    is the volume of the unit $m$-sphere $S^m \subset \R^{m+1}$.
\end{theorem}

\begin{remark}{Why Even Dimension?}{gb-even-dim}
    The formula is false when $M$ is odd-dimensional, since $\chi(M) = 0$ in that case, while $\int_M \kappa$ can be non-zero.
\end{remark}

\begin{proof}
    By the degree formula,
    \[
        \int_M \kappa\, \mathrm{vol}_M = \int_M g^* \mathrm{vol}_{S^m} = \deg(g) \int_{S^m} \mathrm{vol}_{S^m} = \deg(g)\, \gamma_m,
    \]
    so Gauss--Bonnet is equivalent to the claim
    \[
        \chi(M) = 2 \deg(g).
    \]

    This can be shown using the Poincar\'e--Hopf index theorem. Choose $a \in S^m$ such that $\pm a$ are regular values of $g$, and let $v$ be the vector field on $M$ whose value $v_x$ at $x \in M$ is the projection of $-a$ onto $T_x M$.

    Note that $z \in M$ is a zero of $v$ if and only if $g(z) = \pm a$. It follows that $v$ has only finitely many zeros, and since $dv_z = \pm dg_z$ at such a zero,
    \begin{align*}
        \chi(M) = \sum_{v(z) = 0} \mathrm{ind}_z\, v &= \sum_{\substack{v_z = 0 \\ g(z) = a}} \mathrm{ind}_z\, v + \sum_{\substack{v_z = 0 \\ g(z) = -a}} \mathrm{ind}_z\, v \\
        &= \deg(g) + (-1)^m \deg(g) = (1 + (-1)^m) \deg(g) = 2 \deg(g),
    \end{align*}
    using that $m$ is even.
\end{proof}
